% !TEX root = ../main.tex

\documentclass[../main.tex]{subfiles}

\graphicspath{{\subfix{../images/}}}

\begin{document}

В ходе выполнения выпускной квалификационной работы был проведён комплексный анализ предметной области трекинга достижений и геймификации персональной продуктивности. Исследование существующих решений (Habitica, Strides, Any.do, а также специализированных приложений из Google Play) позволило выявить их ключевые преимущества и недостатки: отсутствие стабильного кроссплатформенного продукта с возможностью создания и обмена пользовательскими наборами достижений, ограниченная поддержка пошагового прогресса и зависимость ряда решений от постоянного сетевого подключения. Полученные результаты легли в основу формулировки функциональных и нефункциональных требований к разрабатываемому приложению Achi.

На следующем этапе были проанализированы технологии, позволяющие реализовать приложение в соответствии с сформулированными требованиями. В качестве основы выбран стек Kotlin Multiplatform и Compose Multiplatform, обеспечивающий единую кодовую базу для мобильных платформ (Android, iOS) и Web (WasmJS). Для сетевого взаимодействия применён Ktor Client, для загрузки изображений~--- Coil~3, для навигации~--- библиотека Navigation~3 с типобезопасными маршрутами (NavKey), для генерации API-клиента~--- OpenAPI Generator на основе спецификации \texttt{Achi\_openapi.json}. Архитектура приложения построена по паттерну MVVM с выделением слоя репозиториев; внедрение зависимостей реализовано вручную через контейнер \texttt{AppModule}, что обеспечивает прозрачность связей между компонентами без использования сторонних DI-фреймворков.

В рамках практической реализации разработаны доменные модели данных (пакеты достижений, достижения, шаги с инкрементальным и бинарным прогрессом), репозитории для работы с API и слой представления на базе Compose Multiplatform с ViewModel и реактивным управлением состоянием через \texttt{StateFlow}. Реализованы ключевые пользовательские сценарии: регистрация и аутентификация с сохранением сессии, просмотр коллекции паков, отслеживание прогресса с синхронизацией на сервер, добавление паков по уникальному коду, создание и редактирование собственных паков с загрузкой изображений, локализация интерфейса на русский и английский языки. Для отладочных сборок предусмотрена панель переключения API-хоста.

Финальным этапом работы стало тестирование разработанного решения. Проведены функциональное тестирование основных сценариев использования и кроссплатформенные проверки на целевых платформах (Android, Web и iOS), подтвердившие корректность работы общей логики и пользовательского интерфейса.

Таким образом, все поставленные в работе задачи выполнены. Разработанное приложение Achi удовлетворяет сформулированным требованиям и может быть использовано для отслеживания персональных достижений с элементами геймификации. Перспективными направлениями дальнейшего развития являются реализация полноценного офлайн-режима с локальным хранилищем данных, расширение механизмов обмена паками и добавление in-app переключения языка интерфейса.

\end{document}
